\documentclass[preview]{standalone}

\usepackage{amsmath}
\usepackage{amssymb}
\usepackage{stellar}
\usepackage{definitions}

\begin{document}

\id{quadratic-integers-nonunique-factorization}
\genpage

\section{The ring of integers with square root of minus five}

\begin{snippetexample}{sqrt-minus-five-integer-ring-example}{The ring of integers with square root of minus five}
    Consider the \subring
    \[
        A=\integers[\sqrt{-5}]
        =
        \{a+b\sqrt{-5}\suchthat a,b\in\integers\}
        \subseteq\complexnumbers.
    \]
    Since \(A\) is a \subring of \(\complexnumbers\), it is an
    \integraldomain. Writing \(s=\sqrt{-5}\), one has
    \[
        A=\integers[s],
        \qquad
        s^2=-5.
    \]
\end{snippetexample}

\begin{snippetdefinition}{sqrt-minus-five-norm-definition}{Norm on the integers with square root of minus five}
    For
    \(A=\integers[\sqrt{-5}]\), define the \emph{norm}
    \[
        N\colon A\fromto\naturalnumbers,
        \qquad
        N(a+b\sqrt{-5})=a^2+5b^2.
    \]
    Equivalently, if \(\alpha=a+b\sqrt{-5}\) and
    \(\overline{\alpha}=a-b\sqrt{-5}\), then
    \[
        N(\alpha)=\alpha\overline{\alpha}.
    \]
\end{snippetdefinition}

\begin{snippetproposition}{sqrt-minus-five-norm-multiplicative}{Multiplicativity of the norm}
    For all
    \(\alpha,\beta\in\integers[\sqrt{-5}]\),
    \[
        N(\alpha\beta)=N(\alpha)N(\beta).
    \]
\end{snippetproposition}

\begin{snippetproof}{sqrt-minus-five-norm-multiplicative-proof}{sqrt-minus-five-norm-multiplicative}{Multiplicativity of the norm}
    Since conjugation is multiplicative and the \ring is commutative,
    \begin{align*}
        N(\alpha\beta)
        &=
        \alpha\beta\overline{\alpha\beta}\\
        &=
        \alpha\beta\overline{\alpha}\,\overline{\beta}\\
        &=
        \alpha\overline{\alpha}\,
        \beta\overline{\beta}\\
        &=
        N(\alpha)N(\beta).
    \end{align*}
\end{snippetproof}

\begin{snippetproposition}{sqrt-minus-five-units}{Units in the integers with square root of minus five}
    For \(A=\integers[\sqrt{-5}]\),
    \[
        A^\times=\{\pm1\}.
    \]
\end{snippetproposition}

\begin{snippetproof}{sqrt-minus-five-units-proof}{sqrt-minus-five-units}{Units in the integers with square root of minus five}
    An element \(\alpha\in A\) is a unit \ifandonlyif \(N(\alpha)=1\).
    Indeed, if \(\alpha\beta=1\), multiplicativity gives
    \[
        N(\alpha)N(\beta)=1,
    \]
    so \(N(\alpha)=1\). Conversely, if
    \(\alpha=a+b\sqrt{-5}\) satisfies \(N(\alpha)=1\), then
    \[
        a^2+5b^2=1,
    \]
    forcing \(b=0\) and \(a=\pm1\). Thus \(\alpha=\pm1\), which is a
    unit.
\end{snippetproof}

\section{Irreducibles that are not prime}

\begin{snippet}{sqrt-minus-five-two-factorizations-of-six}
    In \(A=\integers[\sqrt{-5}]\),
    \[
        6=2\cdot3
        =
        (1-\sqrt{-5})(1+\sqrt{-5}).
    \]
    These two decompositions reveal the failure of unique factorization.
\end{snippet}

\begin{snippetproposition}{sqrt-minus-five-irreducible-elements}{Four irreducible elements}
    In \(A=\integers[\sqrt{-5}]\), the elements
    \[
        2,
        \qquad
        3,
        \qquad
        1+\sqrt{-5},
        \qquad
        1-\sqrt{-5}
    \]
    are \irreducibleelement[irreducible elements].
\end{snippetproposition}

\begin{snippetproof}{sqrt-minus-five-irreducible-elements-proof}{sqrt-minus-five-irreducible-elements}{Four irreducible elements}
    If \(2=\alpha\beta\), then
    \[
        4=N(2)=N(\alpha)N(\beta).
    \]
    A nontrivial factorization would require an element of norm \(2\), but
    the equation \(a^2+5b^2=2\) has no integer solutions. Hence one factor
    has norm \(1\), and is a unit.

    Similarly, a nontrivial factorization of \(3\) would require an
    element of norm \(3\), whereas \(a^2+5b^2=3\) has no integer
    solutions. Thus \(3\) is irreducible.

    Finally,
    \[
        N(1+\sqrt{-5})
        =
        N(1-\sqrt{-5})
        =
        6.
    \]
    A nontrivial factorization would require a factor of norm \(2\) or
    \(3\), neither of which exists. Both elements are therefore
    irreducible.
\end{snippetproof}

\begin{snippetproposition}{sqrt-minus-five-irreducible-not-prime}{The four irreducibles are not prime}
    In \(A=\integers[\sqrt{-5}]\), none of
    \[
        2,
        \qquad
        3,
        \qquad
        1+\sqrt{-5},
        \qquad
        1-\sqrt{-5}
    \]
    is a \primeelement.
\end{snippetproposition}

\begin{snippetproof}{sqrt-minus-five-irreducible-not-prime-proof}{sqrt-minus-five-irreducible-not-prime}{The four irreducibles are not prime}
    Since
    \[
        (1-\sqrt{-5})(1+\sqrt{-5})=6\in(2),
    \]
    primality of \((2)\) would require one of
    \(1\pm\sqrt{-5}\) to belong to \((2)\). If, for example,
    \(1+\sqrt{-5}=2\alpha\), then
    \[
        6=N(1+\sqrt{-5})=4N(\alpha),
    \]
    which is impossible. The conjugate case is identical, so \(2\) is
    not prime.

    The same argument for \((3)\) would require
    \(6=9N(\alpha)\), which is impossible. Thus \(3\) is not prime.

    On the other hand,
    \[
        2\cdot3
        =
        (1+\sqrt{-5})(1-\sqrt{-5})
        \in(1+\sqrt{-5}).
    \]
    Neither factor belongs to this principal \ideal: the equations
    \[
        2=(1+\sqrt{-5})\alpha,
        \qquad
        3=(1+\sqrt{-5})\alpha
    \]
    would respectively imply
    \[
        4=6N(\alpha),
        \qquad
        9=6N(\alpha).
    \]
    Both are impossible. Hence \(1+\sqrt{-5}\) is not prime, and
    conjugation gives the same conclusion for \(1-\sqrt{-5}\).
\end{snippetproof}

\begin{snippetcorollary}{sqrt-minus-five-not-principal-ideal-domain}{The integers with square root of minus five are not a PID}
    The \ring \(\integers[\sqrt{-5}]\) is neither a
    \principalidealdomain nor a \euclideandomain.
\end{snippetcorollary}

\begin{snippetproof}{sqrt-minus-five-not-principal-ideal-domain-proof}{sqrt-minus-five-not-principal-ideal-domain}{The integers with square root of minus five are not a PID}
    In every PID, each \irreducibleelement is prime. The preceding
    proposition supplies irreducible elements of
    \(\integers[\sqrt{-5}]\) that are not prime, so this \ring is not a
    PID. Every \euclideandomain is a PID, so it is not Euclidean either.
\end{snippetproof}

\section{Failure of factorization into primes}

\begin{snippetproposition}{sqrt-minus-five-six-no-prime-factorization}{Six has no factorization into prime elements}
    In \(A=\integers[\sqrt{-5}]\), the element \(6\) cannot be written as
    a finite product of \primeelement[prime elements].
\end{snippetproposition}

\begin{snippetproof}{sqrt-minus-five-six-no-prime-factorization-proof}{sqrt-minus-five-six-no-prime-factorization}{Six has no factorization into prime elements}
    Suppose
    \[
        6=p_1\cdots p_r
    \]
    with each \(p_i\) prime. Since \(p_1\divides6=2\cdot3\), primality
    gives \(p_1\divides2\) or \(p_1\divides3\).

    If \(2=p_1\alpha\), irreducibility of \(2\) and the fact that \(p_1\)
    is not a unit imply that \(\alpha\) is a unit. Thus \(p_1\sim2\),
    which would make \(2\) prime, a contradiction. The case
    \(p_1\divides3\) similarly makes \(p_1\sim3\), contradicting that
    \(3\) is not prime. Therefore no such factorization exists.
\end{snippetproof}

\begin{snippetcorollary}{sqrt-minus-five-not-unique-factorization-domain}{The integers with square root of minus five are not a UFD}
    The \ring
    \[
        \integers[\sqrt{-5}]
    \]
    is not a \uniquefactorizationdomain.
\end{snippetcorollary}

\begin{snippetproof}{sqrt-minus-five-not-unique-factorization-domain-proof}{sqrt-minus-five-not-unique-factorization-domain}{The integers with square root of minus five are not a UFD}
    In a UFD every nonzero nonunit factors into prime elements, whereas
    \(6\) has no such factorization in \(\integers[\sqrt{-5}]\).
\end{snippetproof}

\end{document}
